%
% teil1.tex -- Beispiel-File für das Paper
%
% (c) 2020 Prof Dr Andreas Müller, Hochschule Rapperswil
%
% !TEX root = ../../buch.tex
% !TEX encoding = UTF-8
%
\section{Teil 1
\label{zeta:section:teil1}}

\section{Konvergente und divergente Reihen
\label{konvergente-und-divergente-reihen}}

Um die analytische Fortsetzung zu verstehen, muss man sich zuerst mit den Begriffen der konvergenten und divergenten Reihen auskennen.

Unendliche Summen sind Reihen. 
Hierbei gibt es drei Verhaltensweisen, wie sich eine Reihe verhalten kann.

Ist eine reihe konvergent wie zum Beispiel

\[
1 + \frac{1}{2} + \frac{1}{4} + \frac{1}{8} + \frac{1}{16} + \ldots = 2,
\]

so bekommen wir ein endliches Resultat, da die Funktion der Teilsummen irgendwann gegen Null geht.

Ist eine Reihe bestimmt divergent wie zum Beispiel

\[
\sum_{n=1}^{\infty} 1 = 1 + 1 + 1 + 1 + \ldots
\]

oder

\[
1 - \sum_{n=1}^{\infty} 1 = 1 - 1 - 1 - 1 - \ldots,
\]

so geht die Summe entweder gegen Unendlich oder gegen minus Unendlich.

Ist eine Reihe unbestimmt divergent wie zum Beispiel

\[
\sum_{n=0}^{\infty} (-1)^n = 1 - 1 + 1 - 1 + \ldots,
\]

so alterniert das Ergebnis immer und ist somit undefiniert.

Die Summe aller natürlichen Zahlen ist offensichtlich bestimmt divergent. 
Es werden immer grösser werdende Zahlen addiert und somit geht das Ergebnis gegen Unendlich.

